\section{Exact histories, a fixed clique palette, and hardness theorems}
\label{sec:circuit}

Let
\[
G_2=\{X,\CX,\CCX,H\otimes H\}.
\]
Following Rudolph's gate-dependent definition
\cite[Definition~2.2]{rudolph2026universal},\footnote{Definition, theorem,
and appendix numbers for~\cite{rudolph2026universal} refer to the full
version, arXiv:2411.02681v2.} \(\BQPone^{G_2}\) consists of
promise problems decided by a polynomial-time uniform family of exact
\(G_2\) circuits on an all-zero input, with one computational-basis output
bit, perfect completeness, and inverse-polynomial rejection probability in
the no case.  Rudolph's exact error reduction
\cite[Theorem~2.3]{rudolph2026universal} preserves perfect completeness and
stays in any gate set containing \(G_2\); we apply it first so that soundness
is at most \(1/3\).  The superscript is essential: no approximate compilation
or gate-independent identification is used below.

\subsection{Rational histories and the certified palette}

Write \(K=\sqrt2H\).  Replace each \(H\otimes H\) transition by a gain edge
\(K\) on one work bit followed by a loss edge \(K^{-1}=H/\sqrt2\) on the
other.  Their composite is exactly \(H\otimes H\), and each edge is enforced
by two orthonormal rational three-term vectors.  The vectors and their direct
calculation appear in \cref{app:rational-split}.  This construction is due to
Rudolph~\cite[Appendix~D.1]{rudolph2026universal}.

The propagation constraints for \(X,\CX,\CCX\) decompose, on their
constant-size active support, into two-time vectors
\[
(\ket{0,z}-\ket{1,Uz})/\sqrt2.
\]
Clock, input, and output penalties are basis projectors.  Unaffected work
qubits are identities and are implemented by the harmonic padding of
\cref{lem:scaling-padding}; they are not expanded into exponentially many
rank-one terms.

\begin{lemma}[Certified fixed palette]
\label{lem:palette}
Every local constraint in the split \(G_2\) history is implemented by a
member of one fixed family of weighted clique gadgets of total support
locality at most six.  Every member satisfies
\cref{ass:local-certificate,eq:interface}, with uniform constants, and the
graph construction is polynomial in the circuit size.
\end{lemma}

\begin{proof}
See \cref{app:palette}.  It records the finite integer certificates, exact
relabelings, selected-cycle guard closure, and polynomial-size bookkeeping.
\end{proof}

\subsection{Complete history kernels}

Let \(L\) be the number of legal clock positions after splitting the
Hadamard transitions.  Put \(s_t=\sqrt2\) immediately after a gain edge and
\(s_t=1\) otherwise.  If \(V_t\) is the unitary prefix after removing these
scalars, then
\begin{equation}
\label{eq:history-isometry}
\cV v=Z^{-1/2}\sum_{t=0}^{L-1}s_t\ket tV_tv,\qquad
Z=\sum_ts_t^2,\qquad L\le Z\le2L.
\end{equation}
Every gain edge is immediately followed by its matching loss edge, with no
gate between them, so \(s_0=s_{L-1}=1\).

\begin{proposition}[History kernels and gaps]
\label{prop:history}
Let \(C\) be the valid clean-input space, \(D=\dim C>0\), and let
\(M=J^*U^*P_{\rm acc}UJ\) be the compressed acceptance operator.  Put
\(S=\ker(I-M)\).  The initial history Hamiltonian has complete kernel
\(\cV C\), while adding the final rejection projector gives complete kernel
\(\cV S\).  If \(M|_{S^\perp}\preceq I/3\), then
\begin{equation}
\label{eq:history-gaps}
g_{\rm in}\ge\frac1{8L^2},\qquad
g_{\rm out}\ge\frac1{3L(8L^2+1)}\ge\frac1{27L^3}.
\end{equation}
The second statement includes \(S=0\).
\end{proposition}

The proof in \cref{app:source-details} uses a weighted path Poincar\'e
inequality.  The output compression is exactly
\begin{equation}
\label{eq:output-compression}
\cV^*H_{\rm rej}\cV=(I-M)/Z,
\end{equation}
and a two-positive-operator estimate turns this angle penalty into the second
gap in \cref{eq:history-gaps}.  Combining
\cref{thm:transfer}, \cref{thm:quotient}, \cref{lem:palette}, and
\cref{prop:history} yields weighted
clique complexes \(\Xin\subseteq\Xout\) with
\begin{equation}
\label{eq:generic-circuit-transfer}
\beta_d(\Xin)=D,\qquad
\beta_d(\Xout)=\beta_d(\Xin\to\Xout)=\dim S,
\end{equation}
and inverse-polynomial gaps above both complete kernels.

\subsection{A general source: perfect-subspace fractions}
\label{sec:sf}

\Cref{prop:history} places no restriction on the clean-input space \(C\)
beyond \(D>0\), and \cref{eq:generic-circuit-transfer} turns any exact
circuit with a perfectly accepted subspace and a soundness bound on its
complement into a normalized-persistence instance.  We make this source
explicit.

\begin{definition}[Perfect-subspace fraction problems]
\label{def:sf}
Let \(G\) be a finite gate set and let \(a,b:\N\to[0,1]\) satisfy
\(a-b\ge1/\poly\).  A promise problem \(L\) is in \(\SF^{G}(a,b)\) if there
are polynomials \(c,r\) and a polynomial-time uniform family of exact
\(G\)-circuits \(U_x\) acting on \(c(|x|)\) clean qubits, initialized to a
fixed computational-basis string, and \(r(|x|)\) \emph{unconstrained}
qubits, with a computational-basis acceptance projector \(P_{\rm acc}\) on
one output qubit, such that the compressed acceptance operator
\(M_x=J^*U_x^*P_{\rm acc}U_xJ\) on the unconstrained register and its perfect
subspace \(S_x=\ker(I-M_x)\) satisfy
\begin{enumerate}[(i)]
\item \(M_x|_{S_x^\perp}\preceq\tfrac13 I\), that is, every state orthogonal
to the perfectly accepted subspace is accepted with probability at most
\(1/3\);
\item \(\dim S_x/2^{r}\ge a(|x|)\) if \(x\in L_{\rm yes}\), and
\(\dim S_x/2^{r}\le b(|x|)\) if \(x\in L_{\rm no}\).
\end{enumerate}
\end{definition}

Three sources of this type matter here.

\begin{itemize}
\item \emph{Perfect-completeness circuits.}  \(\BQPone^{G_2}\subseteq\SF^{G_2}(3/4,1/8)\)
with \(r=3\): given a \(\BQPone^{G_2}\) verifier \(Q_x\) with acceptance
probability \(p_x\), add a three-bit unconstrained label \(z\), run \(Q_x\)
unconditionally, and accept according to
\[
[z=000]\ \lor\
\bigl(q\land[z\in\{001,010,011,100,101\}]\bigr),
\]
where \(q\) is the output of \(Q_x\).  The reversible predicate uses only
\(X,\CX,\CCX\), preserves \(z\), and uncomputes its scratch.  Its compressed
acceptance operator on the label is exactly
\begin{equation}
\label{eq:eight-label}
M_x=\operatorname{diag}(1,p_x,p_x,p_x,p_x,p_x,0,0),
\end{equation}
so the perfect subspace has dimension six when \(p_x=1\) and dimension one
when \(p_x\le1/3\), and every off-perfect eigenvalue in the latter case is at
most \(1/3\) (\cref{app:source-boundaries}).
\item \emph{One clean qubit.}  Lowe et al.\ define
\(\SDQCone^{G}\)~\cite[Definition~2]{lowe2026normalized} by a uniform
\(G\)-circuit on one clean and \(q-1\) maximally mixed qubits, a subspace
\(S_x\) of the mixed register accepted with probability exactly one, soundness
\(1/3\) on \(S_x^\perp\), and a decision condition on the acceptance
probability of the maximally mixed input.  Conditions (i)--(ii) of their
definition give \(S_x=\ker(I-M_x)\) and \(M_x|_{S_x^\perp}\preceq I/3\)
(\cref{app:source-boundaries}).  If the decision condition is imposed on the
fraction \(\dim S_x/2^{q-1}\), which is the quantity that their Section~3.1
identifies as the power of the class and from which the instance is decided
in the proofs of their Theorems~5 and~7, the resulting class, which we write
\(\SDQCone^{G}[\mathrm{frac}]\), is contained in \(\SF^{G}(a,b)\) with \(c=1\).
With the decision condition on the acceptance probability \(p_x\) of the
maximally mixed state, the identity
\(f\le p_x\le(1+2f)/3\) for \(f=\dim S_x/2^{q-1}\)
(\cref{app:source-boundaries}) gives
\(\SDQCone^{G}(a,b)\subseteq\SF^{G}\bigl(\tfrac{3a-1}{2},b\bigr)\) whenever
\(b<(3a-1)/2\).
\item \emph{Classical counting.}  A Boolean formula \(\varphi\) on \(r\)
variables, evaluated by a reversible \(\{X,\CX,\CCX\}\) circuit on the
unconstrained register with clean scratch, has \(M_\varphi\) diagonal with
entries \(\varphi(z)\in\{0,1\}\); so \(S_\varphi\) is spanned by the satisfying
assignments, \(\dim S_\varphi=\#\varphi\), and (i) holds with
\(M_\varphi|_{S^\perp}=0\).
\end{itemize}

\begin{theorem}[General hardness]
\label{thm:general-hardness}
For every \(a-b\ge1/\poly\) there is a deterministic polynomial-time many-one
reduction from \(\SF^{G_2}(a,b)\) to weighted gapped normalized persistence
(\cref{def:problem}) with additive error \(\varepsilon<(a-b)/2\).  On input
\(x\) it outputs \(\Xin\subseteq\Xout\) with
\[
\beta_d(\Xin)=2^{r},\quad
\beta_d(\Xin\to\Xout)=\beta_d(\Xout)=\dim S_x,\quad
q_d(\Xin,\Xout)=\frac{\dim S_x}{2^{r}},
\]
and explicit rational bounds \(E_{\rm in},E_{\rm out}\ge1/\poly\) below the
positive spectra of both endpoint Laplacians.
\end{theorem}

\begin{proof}
Take \(C\) to be the fixed clean string tensored with the full unconstrained
register, so \(D=2^{r}\), and let \(S=S_x\).  Condition (i) is the hypothesis
\(M|_{S^\perp}\preceq I/3\) of \cref{prop:history}, so both history kernels
and gaps of \cref{eq:history-gaps} are available.  Apply
\cref{eq:generic-circuit-transfer}.  The YES and NO values of \(q_d\) are
separated by \(a-b\), so additive error below \((a-b)/2\) decides.  Parameter
and size bookkeeping: the logical gaps in \cref{eq:history-gaps} are explicit
rationals, and the fixed palette supplies hard-coded positive rational lower
bounds for the constants in \cref{thm:transfer}.  Fix \(\eta_0=1/10\); from
the polynomially bounded \(t_{\max}\), choose the first integer \(b'\) for
which \(\lambda=2^{-b'}\) satisfies both inequalities in
\cref{eq:lambda-E-choice}, and for each endpoint output the largest dyadic
rational \(E_i\) no greater than one half of \(c_2\eta_0^2g_i\lambda^{2p}\).
Thus \(b'=O(\log t_{\max})\) and the binary encodings of \(\lambda\) and
\(E_i\) have polynomial length.  If the history uses \(n\) register qubits
(clock and work bits), the register has \(7n\) vertices and \(O(n^2)\)
edges, and each gadget adds a constant number of private vertices together
with \(O(n)\) edges to the unaffected register factors; so the weighted
graph has \(O(n+t_{\max})\) vertices and \(O(n^2+t_{\max}n)\) explicitly
generated edges.  Hence the map outputs the two graphs, their common weight
function, \(d\), \(E_{\rm in},E_{\rm out}\), and the threshold
\((a+b)/2\) in deterministic polynomial time.
\end{proof}

\begin{corollary}[One-clean-qubit sources]
\label{cor:sdqc}
Weighted gapped normalized persistence is \(\SDQCone^{G_2}[\mathrm{frac}]\)-hard
for all thresholds \(a-b\ge1/\poly\), and \(\SDQCone^{G_2}\)-hard, in the
sense of \cite[Definition~2]{lowe2026normalized}, for all thresholds with
\(b<(3a-1)/2\); for \((a,b)=(2/3,1/3)\) the target error \(1/24\) suffices.
In these instances \(\beta_d(\Xin)=2^{q-1}\) is the full dimension of the
mixed register.
\end{corollary}

Conjecture~1 of Lowe et al.\ asserts \(\SDQCone\)-hardness of this problem
for a fixed universal gate set.  \Cref{cor:sdqc} proves it for \(G_2\) in the
fraction form used in their own hardness proofs, and in the trace form for
every threshold pair with \(b<(3a-1)/2\).  Arbitrary trace thresholds
cannot be handled by any reduction whose target sees only the fraction
\(\dim S_x/2^{q-1}\) (\cref{app:source-boundaries}).

\begin{theorem}[Fixed-eight hardness]
\label{thm:weighted-hardness}
There is a deterministic polynomial-time many-one reduction from
\(\BQPone^{G_2}\) to the decision version of weighted gapped normalized
persistence (\cref{def:problem}).  Every output
satisfies
\[
\beta_d(\Xin)=8,\qquad
\beta_d(\Xin\to\Xout)=
\begin{cases}6,&x\in L,\\1,&x\notin L,\end{cases}
\]
and both endpoint degree-\(d\) Laplacians have inverse-polynomial gaps above
their complete kernels.  Hence the normalized values are \(3/4\) and \(1/8\).
Additive error \(1/24\), with decision threshold \(7/16\), distinguishes the
two cases with success probability at least \(2/3\).
\end{theorem}

\begin{proof}
Apply \cref{thm:general-hardness} to the eight-label source
\cref{eq:eight-label}, with \(r=3\).  The two ratios differ by \(5/8\), and
\(1/24<5/16\).
\end{proof}

\begin{theorem}[Gapped Betti numbers]
\label{thm:gapped-betti}
Consider the problem \textnormal{\textsc{GappedBetti}}: given a graph \(G\) with dyadic
vertex weights, a degree \(d\), and a rational \(\gamma\ge1/\poly\), with
the promise \(\spec(\Delta_{\Cl(G),d})\subseteq\{0\}\cup[\gamma,\infty)\),
compute \(\beta_d(\Cl(G))\).  \textnormal{\textsc{GappedBetti}} is \(\#\mathsf P\)-hard under
polynomial-time parsimonious reductions, and it is in
\(\#\BQP\)~\cite{brown2011computational}, hence \(\#\mathsf P\)-complete under
weakly parsimonious reductions.  The same statements hold for the persistent
Betti number \(\beta_d(X_1\to X_2)\) of a filtration with gapped endpoints,
and with exact \(G_2\) verifiers satisfying condition (i) of \cref{def:sf} in
place of Boolean formulas, which gives hardness, again parsimoniously, for
the perfect-completeness counting class \(\#\BQPone\) of Cade and
Crichigno~\cite{cade2024complexity} restricted to \(G_2\).
\end{theorem}

\begin{proof}
Hardness: given \(\varphi\), the classical-counting source above satisfies
(i), and \cref{thm:general-hardness}, which uses only the two-term atoms and
basis cones of the palette for \(X,\CX,\CCX\) circuits, outputs \(\Xout\)
with \(\beta_d(\Xout)=\beta_d(\Xin\to\Xout)=\#\varphi\), an inverse-polynomial
gap above the kernel, and \(\beta_d(\Xin)=2^r\).  Containment: with the
promised gap, phase estimation of the explicit sparse operator
\(\Delta_{\Cl(G),d}\) at resolution below \(\gamma\) is a \(\BQP\) verifier
that accepts kernel vectors with probability exponentially close to one and
every vector of the orthogonal complement with probability exponentially
close to zero, so \(\beta_d\) is the dimension of its accepting subspace, a
\(\#\BQP\) function~\cite{brown2011computational,cade2024complexity}; and
\(\#\BQP=\#\mathsf P\) under weakly parsimonious
reductions~\cite{brown2011computational}.
\end{proof}

Crichigno and Kohler prove \(\#\mathsf P\)-hardness of Betti numbers of clique
complexes without a gap promise~\cite{crichigno2024clique}, and Cade and
Crichigno prove the gapped statement for general chain complexes with
\(6\)-local coboundaries~\cite{cade2024complexity}; \cref{thm:gapped-betti}
combines the two features for clique complexes, and, through
\cref{cor:unweighted-hardness} below, for unweighted clique complexes.

The exact trace-versus-perfect-space obstruction and the mixed-spectator
extension to individual real Hadamards are recorded in
\cref{app:source-boundaries}.  Neither is needed for
\cref{thm:weighted-hardness}.
